(1) $P_k^2(a) = P_k\bigl((a\cdot k)k\bigr) = (a\cdot k)\,(Q(k))\,k = 0$ since $k^2 = 0$; the group law and the closed-form inverse follow exactly as in Proposition \claimref{interior.gauge-trajectory-unipotent-ray}, with parameter $s = H$ varying over the manifold. (2) Expand: $\hb(a)\cdot\hb(b) = a\cdot b + H(b\cdot k)(a\cdot k) + H(a\cdot k)(k\cdot b) + H^2 (a\cdot k)(b\cdot k)\,k^2$, and the last term vanishes; the middle terms add. The inverse case is the same expansion with $H \mapsto -H$, using $\hb^{-1} = 1 - HP_k$ from (1); the sign of the metric correction is the one the paper's definition $g(a,b) = \hb^{-1}(a)\cdot\hb^{-1}(b)$ forces, and at $a = 0$, $H = M/r$, $k = \et - \er$ it reproduces ingoing Eddington-Finkelstein Schwarzschild with the correct mass sign (the metric dual of $\et - \er$ is $dt + dr$, and $\eta - 2H(dt + dr)^2$ is the ingoing form in the advanced time $v = t + r$; the outgoing congruence is $\et + \er$). (3) is (1) and (2) applied to Fact \claimref{interior.imported-kerr-schild-form-kerr-gauge}, plus the comparison with \eqref{eq:hbar}: there $P(a) = (a\cdot\et)\er$ with $\et\cdot\er = 0$ but $\et, \er$ non-null; here the single vector $k$ is null, so the map $P_k$ shears along the very direction it projects onto. The parabolic reading of Remark \ref{rem:parabolic} therefore applies to the whole family with one refinement: rotation moves the gauge trajectory onto the null boundary of the shear cone from the start.
